% !TeX root = ./../main.tex

\section{Próbkowanie kwantowe jako model obliczeń}\label{sec:tb:qsamp}
\agentinfo{
Stan (PHD-TS-95, 2026-11-27): podrozdział skrócony i poprawiony po recenzji \agentref{eksperymenty/recenzje/rozdzial-2-sekcje-2.3-2.4.md}; usunięto podrozdział o operatorze zniekształcenia fazy (definicja, twierdzenie o odczycie), funkcję FCoSamp z postacią Fouriera i powtórzony wzór interferencji; lista cięć w \agentref{eksperymenty/recenzje/rozdzial-2-skracanie.md}.
QCoSamp i kodowanie danych w fazie pochodzą z pracy Wereszczyńskiego i in. \agentref{bib:wereszczynski_cosine_2020} (preprint) i nie są wkładem autora rozprawy.
Konwencja (za pracą źródłową i kodem z pracy magisterskiej \agentref{artykuły/magisterka_werner_final/chapters/03.tex}): argument x_D = -pi + suma po j od 0 do m-1 z pi d_j / 2^j, bit d_0 najbardziej znaczący, siatka od -pi do pi - pi/2^(m-1); wartość składnika jest prawdopodobieństwem wyniku 0 na kubicie odczytu.
Rozbieżności źródeł niewidoczne po skróceniu: warunek lambda^2 + gamma^2 <= 2 w pracy źródłowej (z tożsamości trygonometrycznej wynika <= 4) i wzór odtwarzania szeregu Fouriera dotyczyły usuniętego akapitu o funkcji FCoSamp; dziedzina [-1, 1] funkcji referencyjnej w artykule ICCS \agentref{artykuły/Quantum_Errors_ICCS/main.tex} należy do rozdziału 4.
Etykiety używane przez rozdział 4 (nie usuwać): subsec:tb:qsamp:sampling, subsec:tb:qsamp:qcosamp, def:tb:qsamp:sampling, def:tb:qsamp:component, def:tb:qsamp:argument, def:tb:qsamp:circuit, thm:tb:qsamp:readout, fig:tb:qsamp:component, eq:tb:qsamp:two-variables.
}

Metody redukcji błędów rozwijane w rozprawie korygują wynik obliczeń, który ma postać rozkładu prawdopodobieństwa, a nie pojedynczego ciągu bitów.
Podrozdział definiuje ten sposób prowadzenia obliczeń i metodę QCoSamp, na obwodach której opracowano metodę PDU (rozdział~\ref{chapter:pdu}); metoda ta i kodowanie danych w fazie pochodzą z pracy Wereszczyńskiego i~in.~\cite{wereszczynski_cosine_2020}, dostępnej jako preprint, i nie są wkładem autora rozprawy.

\subsection{Obliczenia oparte na próbkowaniu}\label{subsec:tb:qsamp:sampling}

Pojedyncze wykonanie obwodu kończy się pomiarem, który zwraca jeden stan bazowy (podrozdział~\ref{subsec:tb:measurement}); w \glslink{quantum-sampling}{próbkowaniu kwantowym} wynikiem obliczeń nie jest pojedynczy taki stan, lecz rozkład prawdopodobieństwa stanów bazowych, szacowany z wielu wykonań tego samego obwodu~\cite{wereszczynski_cosine_2020,werner_experiment_driven_2022}.

\begin{Definition}[Eksperyment próbkowania i znormalizowany histogram]\label{def:tb:qsamp:sampling}
Niech $U$ będzie operatorem ewolucji obwodu, a $\mathbb{E}$ bazą pomiarową złożoną ze stanów bazowych mierzonych kubitów.
\emph{Eksperymentem próbkowania} o liczbie powtórzeń $\nshots$ nazywa się $\nshots$-krotne przygotowanie stanu $U\ket{0\dots0}$ i jego pomiar w bazie $\mathbb{E}$, a \emph{znormalizowanym histogramem} eksperymentu funkcję $\tilde{\prob}(\ket{k})$ określoną na $\mathbb{E}$, której wartościami są częstości względne~\eqref{eq:tb01:estimator} wyników.
\end{Definition}

Oprócz obwodu i eksperymentu próbkowania obliczenie wymaga słownika, który stanom bazowym przypisuje znaczenie w dziedzinie problemu (wartość funkcji w zadanym punkcie, intensywność piksela, prawdopodobieństwo przynależności obiektu do klasy), oraz rekonstrukcji tych wielkości z histogramu~\cite{wereszczynski_cosine_2020}.

\subsection{Metoda QCoSamp}\label{subsec:tb:qsamp:qcosamp}

Metoda \gls{qcosamp}, czyli kwantowe próbkowanie szeregiem cosinusów, służy do budowy operatorów ewolucji, dla których prawdopodobieństwo wyniku pomiaru jednego kubitu jest ważoną sumą przesuniętych funkcji cosinus argumentu wprowadzonego do obwodu~\cite{wereszczynski_cosine_2020}.

\begin{Definition}[Składnik cosinusowy]\label{def:tb:qsamp:component}
\emph{Składnikiem cosinusowym} o częstotliwości (krotności argumentu) $n\in\{1,2,\dots\}$ i fazie $r\in\mathbb{R}$ nazywa się funkcję o wartościach w przedziale $[0,1]$
\begin{equation}\label{eq:tb:qsamp:component}
\coscomp_{n,r}(x)=\frac{1}{2}\bigl(1+\cos(nx+r)\bigr),\qquad x\in[-\pi,\pi).
\end{equation}
\end{Definition}

Wartość składnika powstaje przez \glslink{interference}{interferencję} na jednym kubicie: zgodnie z~\eqref{eq:tb01:interference} układ bramek $\mathrm{H}\,\mathrm{P}(\varphi)\,\mathrm{H}$ daje $\prob(\ket{0})=\coscomp_{n,r}(x)$ dla $\varphi=nx+r$, a argument $x$ wprowadza się do obwodu jako stan bazowy osobnego rejestru.
W pracy źródłowej składnikiem (ang.\ \emph{component}) nazywa się parę takich funkcji o wspólnej częstotliwości, a kubit, na którym zachodzi interferencja, kubitem pomocniczym (\emph{ancilla}); w rozprawie składnikiem jest pojedyncza funkcja, a kubit ten nazywa się kubitem odczytu.
\agenttask{
Pytanie do autora (recenzja 2.4-I1): czy zaakceptować odstępstwa od zapisu pracy źródłowej metody QCoSamp (składnik jako pojedyncza funkcja, nazwa kubit odczytu zamiast ancilla, oznaczenia własne składnika, obwodu i bloku fazowego) i czy potwierdzić u współautorów preprintu dwie usterki jego zapisu, których rozprawa po skróceniu już nie przytacza (warunek na współczynniki z prawą stroną 2 zamiast 4 oraz wzór odtwarzania bez odjęcia wyrazu stałego)?
}

\begin{Definition}[Kodowanie argumentu]\label{def:tb:qsamp:argument}
Niech $D=(d_0,\dots,d_{m-1})\in\{0,1\}^{m}$ będzie ciągiem bitów zapisanym w stanie bazowym $\ket{D}=\ket{d_0\dots d_{m-1}}$ \emph{rejestru argumentu} o $m$ kubitach.
Ciąg $D$ koduje argument
\begin{equation}\label{eq:tb:qsamp:argument}
x_D=-\pi+\sum_{j=0}^{m-1}\frac{\pi d_j}{2^{j}} ,
\end{equation}
a liczbę $2^{m}$ kodowanych wartości nazywa się \emph{rozdzielczością} argumentu.
\end{Definition}

Bit $d_0$ jest najbardziej znaczący, a kodowane argumenty tworzą siatkę równomierną o kroku $\Delta x=\pi/2^{m-1}$ od $-\pi$ do $\pi-\Delta x$.

\begin{Definition}[Obwód składnika]\label{def:tb:qsamp:circuit}
\emph{Obwodem składnika} $\coscomp_{n,r}$ nazywa się obwód o $m+1$ kubitach, złożony z rejestru argumentu i \emph{kubitu odczytu} $a$, jedynego kubitu podlegającego pomiarowi, inicjowany stanem $\ket{0}^{\otimes(m+1)}$ i opisany operatorem
\begin{equation}\label{eq:tb:qsamp:circuit}
\compcirc_{n,r}(D)=\mathrm{H}_a\,\compphase^{(a)}_{n,r}\,\mathrm{H}_a\,E(D),
\quad
\compphase^{(a)}_{n,r}=\Bigl[\,\prod_{j=0}^{m-1}\mathrm{CP}_{j\to a}\bigl(\tfrac{n\pi}{2^{j}}\bigr)\Bigr]\mathrm{P}_a(r-n\pi),
\quad
E(D)=\bigotimes_{j=0}^{m-1}\mathrm{X}^{d_j},
\end{equation}
gdzie $\mathrm{CP}_{j\to a}(\theta)$ oznacza bramkę przesunięcia fazy o kąt $\theta$ na kubicie $a$ kontrolowaną $j$-tym kubitem rejestru argumentu.
\end{Definition}

\begin{Theorem}[Odczyt składnika]\label{thm:tb:qsamp:readout}
Po ewolucji $\compcirc_{n,r}(D)\ket{0}^{\otimes(m+1)}$ prawdopodobieństwa wyników pomiaru kubitu odczytu wynoszą
\begin{equation}\label{eq:tb:qsamp:readout}
\prob(\ket{0}_a)=\coscomp_{n,r}(x_D),\qquad \prob(\ket{1}_a)=1-\coscomp_{n,r}(x_D).
\end{equation}
\end{Theorem}
\begin{proof}
Operator $E(D)$ przeprowadza rejestr argumentu w stan $\ket{D}$, a diagonalny blok $\compphase^{(a)}_{n,r}$ nie zmienia składowej $\ket{D}\ket{0}_a$ i zwiększa fazę składowej $\ket{D}\ket{1}_a$ o $r-n\pi+\sum_{j}n\pi d_j/2^{j}$, czyli o $nx_D+r$, więc teza wynika z~\eqref{eq:tb01:interference} dla $\varphi=nx_D+r$.
\end{proof}

\begin{figure}[htbp]
\centering
\begin{subfigure}[b]{0.58\textwidth}
\centering
% !TeX root = ./../../main.tex
% Obwód składnika cosinusowego QCoSamp dla rejestru argumentu o m = 4 kubitach.
% Rysunek własny (quantikz2); układ bramek zgodny z obwodem referencyjnym
% z pracy magisterskiej autora (artykuły/magisterka_werner_final/chapters/03.tex, listing lst:python_cos_phi).
% Obwód jest składany w pudełku i skalowany do szerokości otoczenia (\linewidth).
\ifdefined\tbqsampbox\else\newsavebox{\tbqsampbox}\fi
\begin{lrbox}{\tbqsampbox}
\begin{quantikz}[row sep=0.14cm, column sep=0.18cm, font=\small]
\lstick{$\ket{0}$} & \gate{\mathrm{X}^{d_0}}\gategroup[4,steps=1,style={dashed,rounded corners,inner sep=2pt},label style={yshift=1pt}]{$E(D)$} & & \gategroup[5,steps=5,style={dashed,rounded corners,inner sep=2pt},label style={yshift=1pt}]{$\compphase^{(a)}_{n,r}$} & \ctrl{4} & & & & & \\
\lstick{$\ket{0}$} & \gate{\mathrm{X}^{d_1}} & & & & \ctrl{3} & & & & \\
\lstick{$\ket{0}$} & \gate{\mathrm{X}^{d_2}} & & & & & \ctrl{2} & & & \\
\lstick{$\ket{0}$} & \gate{\mathrm{X}^{d_3}} & & & & & & \ctrl{1} & & \\
\lstick{$a\colon\ket{0}$} & & \gate{\mathrm{H}} & \gate{\mathrm{P}(r-n\pi)} & \gate{\mathrm{P}(n\pi)} & \gate{\mathrm{P}\bigl(\tfrac{n\pi}{2}\bigr)} & \gate{\mathrm{P}\bigl(\tfrac{n\pi}{4}\bigr)} & \gate{\mathrm{P}\bigl(\tfrac{n\pi}{8}\bigr)} & \gate{\mathrm{H}} & \meter{}
\end{quantikz}
\end{lrbox}
\resizebox{\linewidth}{!}{\usebox{\tbqsampbox}}

\caption{składnik $\coscomp_{n,r}$ dla $m=4$}
\label{fig:tb:qsamp:component}
\end{subfigure}\hfill
\begin{subfigure}[b]{0.39\textwidth}
\centering
% !TeX root = ./../../main.tex
% Schemat blokowy obwodu realizującego szereg dwuskładnikowy (1/2)(kappa_{1,0} + kappa_{2,0}).
% Rysunek własny (quantikz2); struktura zgodna z obwodem z pracy magisterskiej autora
% (artykuły/magisterka_werner_final/graf/Eksperymenty/wiele_komp/uklad.png).
% Obwód jest składany w pudełku i skalowany do szerokości otoczenia (\linewidth).
\ifdefined\tbqsampbox\else\newsavebox{\tbqsampbox}\fi
\begin{lrbox}{\tbqsampbox}
\begin{quantikz}[row sep=0.2cm, column sep=0.2cm, font=\small]
\lstick{$b\colon\ket{0}$} & \gate{\mathrm{H}} & \gate[2]{\compphase^{(b)}_{2,0}} & & \gate[3]{\bigl(\compphase^{(ab)}_{3,0}\bigr)^{\dagger}} & \ctrl{2} & & \\
\lstick{$\ket{D}$} & \qwbundle{m} & & \gate[2]{\compphase^{(a)}_{1,0}} & & & & \\
\lstick{$a\colon\ket{0}$} & \gate{\mathrm{H}} & & & & \targ{} & \gate{\mathrm{H}} & \meter{}
\end{quantikz}
\end{lrbox}
\resizebox{\linewidth}{!}{\usebox{\tbqsampbox}}

\caption{szereg $\tfrac{1}{2}(\coscomp_{1,0}+\coscomp_{2,0})$}
\label{fig:tb:qsamp:two-components}
\end{subfigure}
\caption{Obwody metody QCoSamp: (a)~obwód składnika cosinusowego (definicja~\ref{def:tb:qsamp:circuit}) z kątami bramek fazowych dla $n=1$ i $r=0$; (b)~schemat obwodu szeregu dwuskładnikowego, w którym blok $(\compphase^{(ab)}_{3,0})^{\dagger}$, dodatkowo kontrolowany kubitem $a$, usuwa iloczyn faz $\mathrm{e}^{\mathrm{i}3x_D}$ ze składowej $\ket{1}_a\ket{1}_b$.
Opracowanie własne na podstawie~\cite{wereszczynski_cosine_2020,werner_magisterka_2022}.}
\label{fig:tb:qsamp:circuits}
\end{figure}

Obwód składnika (rysunek~\ref{fig:tb:qsamp:component}) wymaga $m$ kontrolowanych bramek fazowych, a stablicowanie składnika na całej siatce $2^{m}$ eksperymentów próbkowania, po jednym dla każdego ciągu $D$.
W odróżnieniu od pojedynczej bramki $\mathrm{P}(nx+r)$, która daje tę samą wartość, angażuje on wiele kubitów i \glslink{phase-kickback}{przeniesienie fazy} z bramek dwukubitowych (podrozdział~\ref{subsec:tb:gates}), dlatego to on służy w rozprawie do badania błędów.
Dla $n=1$ i $r=0$ otrzymuje się funkcję $\coscomp_{1,0}(x)=\tfrac{1}{2}(1+\cos x)$ i obwód $\compcirc_{1,0}(D)$, które w metodzie PDU pełnią rolę funkcji referencyjnej i obwodu referencyjnego (podrozdział~\ref{sec:pdu:reference}).

\emph{Szeregiem wieloskładnikowym} nazywa się kombinację wypukłą składników cosinusowych, $\cosseries(x)=\sum_{c}w_c\,\coscomp_{n_c,r_c}(x)$, gdzie $w_c>0$ i $\sum_{c}w_c=1$.
W obwodzie realizującym szereg składnik wskazują stany bazowe $\ket{c}$ dodatkowych kubitów pomocniczych, wprowadzonych w superpozycję: jeżeli przed końcową bramką Hadamarda układ jest w stanie $\sum_{c}\sqrt{w_c}\,\ket{c}\otimes(\ket{0}+\mathrm{e}^{\mathrm{i}\varphi_c}\ket{1})_a/\sqrt{2}$, to ze względu na ortogonalność stanów $\ket{c}$ składniki nie interferują ze sobą i $\prob(\ket{0}_a)=\sum_{c}w_c(1+\cos\varphi_c)/2$, czyli $\cosseries(x_D)$ dla $\varphi_c=n_cx_D+r_c$; niepożądane iloczyny faz różnych składników usuwa się przez \glslink{uncomputation}{dekomputację}.
W obwodzie z rysunku~\ref{fig:tb:qsamp:two-components}, użytym w pracy magisterskiej autora~\cite{werner_magisterka_2022}, kubit odczytu $a$ otrzymuje fazę $x_D$, kubit pomocniczy $b$ fazę $2x_D$, a po dekomputacji stan kubitu $b$ wskazuje, którą z~dwóch faz niesie kubit odczytu, więc $\prob(\ket{0}_a)=\tfrac{1}{4}(1+\cos x_D)+\tfrac{1}{4}(1+\cos 2x_D)$.

Funkcję dwóch zmiennych otrzymuje się, dodając drugi rejestr argumentu, przechowujący ciąg $D'$, i kontrolując nim drugi blok fazowy, o częstotliwości $n'$ i zerowej fazie, na tym samym kubicie odczytu; fazy pochodzące od obu rejestrów sumują się i
\begin{equation}\label{eq:tb:qsamp:two-variables}
\prob(\ket{0}_a)=\frac{1}{2}\bigl(1+\cos(nx_D+n'y_{D'}+r)\bigr)=\coscomp_{n,\,r+n'y_{D'}}(x_D),
\end{equation}
gdzie $y_{D'}$ jest argumentem kodowanym przez $D'$ zgodnie z~\eqref{eq:tb:qsamp:argument}.
Konstrukcję tego typu autor zastosował do funkcji zapisanej w pracy~\cite{werner_magisterka_2022} jako $\tfrac{1}{2}(1+\cos(x+y))$ (podrozdział~\ref{subsec:pdu:results:two}); stałego przesunięcia fazy w tamtym obwodzie źródło nie podaje jednoznacznie.

Drugim sposobem wprowadzania danych, opisanym w pracy źródłowej jako kodowanie danych stałych, jest zapisanie wektora liczb w \glslink{local-phase}{fazach lokalnych} stanów bazowych operatorem diagonalnym~\cite{wereszczynski_cosine_2020}; opiera się na nim reprezentacja LPIQE, zdefiniowana wraz z odczytem tych faz w podrozdziale~\ref{subsec:tb:qimg:lpiqe}.
\agenttask{
Pytanie do autora (recenzja 2.4-I3): po skróceniu rozprawa nie używa nazwy operator zniekształcenia fazy (PDO), która pochodzi z koncepcji z 2022 r., a nie ze źródeł opublikowanych, i myliła się z nazwą PDU; czy nazwa ma wrócić (jedno zdanie w tym akapicie albo w definicji LPIQE)?
}
